\documentclass[a4paper,12pt]{article}
\usepackage[utf8]{inputenc}
\usepackage[T1]{fontenc}
\usepackage[ngerman]{babel}
\usepackage{amsmath}
\usepackage{amssymb}
\usepackage{enumitem}
\usepackage{geometry}
\geometry{a4paper, margin=2.5cm}

\title{Einsendeaufgaben -- Lektion 2\\[0.5em]
\large Modul 61111: Mathematische Grundlagen}
\author{}
\date{}

\begin{document}
\maketitle

\section*{Aufgabe 2.5}

Sei $A$ die Matrix mit den gegebenen Vektoren als Spalten:
\[
A = \begin{pmatrix}1 & 1 & 1 \\ 2t & t & 2 \\ 0 & t+2 & 2t\end{pmatrix}.
\]

\noindent\textbf{Fall $t \neq 0$:}

\[
A
\xrightarrow{Z_{2,1}(-2t)}
\begin{pmatrix}1 & 1 & 1 \\ 0 & -t & 2-2t \\ 0 & t+2 & 2t\end{pmatrix}
\xrightarrow{Z_{3,2}(1)}
\begin{pmatrix}1 & 1 & 1 \\ 0 & -t & 2-2t \\ 0 & 2 & 2\end{pmatrix}
\]

\[
\xrightarrow{Z_2\!\left(-\tfrac{1}{t}\right),\; Z_3\!\left(\tfrac{1}{2}\right)}
\begin{pmatrix}1 & 1 & 1 \\ 0 & 1 & \tfrac{2}{t}-2 \\ 0 & 1 & 1\end{pmatrix}
\xrightarrow{Z_{3,2}(-1)}
\begin{pmatrix}1 & 1 & 1 \\ 0 & 1 & \tfrac{2}{t}-2 \\ 0 & 0 & 3-\tfrac{2}{t}\end{pmatrix}
\]

\[
\xrightarrow{Z_{1,2}(-1),\; Z_{2,3}\!\left(-\tfrac{2}{t}+2\right),\; Z_3\!\left(\tfrac{t}{3t-2}\right)}
\begin{pmatrix}1 & 0 & 0 \\ 0 & 1 & 0 \\ 0 & 0 & 1\end{pmatrix}
\]

Die Matrix besitzt 3 Pivot-Positionen. Folglich erzeugt die Teilmenge für alle
$t \in \mathbb{R} \setminus \{0\}$ ganz $\mathbb{R}^3$.

\bigskip
\noindent\textbf{Fall $t = 0$:}

\[
A\big|_{t=0} = \begin{pmatrix}1 & 1 & 1 \\ 0 & 0 & 2 \\ 0 & 2 & 0\end{pmatrix}
\xrightarrow{\text{Zeilentausch}}
\begin{pmatrix}1 & 1 & 1 \\ 0 & 2 & 0 \\ 0 & 0 & 2\end{pmatrix}.
\]

Auch hier besitzt die Matrix 3 Pivot-Positionen, die Teilmenge erzeugt also ganz $\mathbb{R}^3$.

\bigskip
Also gilt für alle $t \in \mathbb{R}$, dass die Teilmenge ganz $\mathbb{R}^3$ erzeugt.

\end{document}
